\subsection{Decomposition group and inertia group}

DEFINITION 2. Let A' be a ring and G a group operating on A' (§ 1, no. 9). Given a prime ideal p' of A' the subgroup of elements σ ∈ G such that σ.p' = p' is called the decomposition group of p' (with respect to G) and is denoted by \(G^{\mathrm{Z}}(p')\). The ring of elements of A' invariant under \(G^{\mathrm{Z}}(p')\) is called the decomposition ring of p' (with respect to G) and is denoted by \(A^{\mathrm{Z}}(p')\) (*).

  We often write \(\mathcal{G}^{\mathbf{Z}}\) and \(\mathbf{A}^{\mathbf{Z}}\) instead of \(\mathcal{G}^{\mathbf{Z}}(\mathfrak{p}')\) and \(\mathbf{A}^{\mathbf{Z}}(\mathfrak{p}')\) respectively, when there is no ambiguity.

For all \(\sigma \in \mathcal{G}^{\mathbb{Z}}(\mathfrak{p}')\) we also denote by \(z \mapsto \sigma.z\) the endomorphism of the ring \(\mathbf{A}' / \mathfrak{p}'\) derived from the endomorphism \(x \mapsto \sigma.x\) of \(\mathbf{A}'\) by taking quotients; clearly the group \(\mathcal{G}^{\mathbb{Z}}(\mathfrak{p}')\) operates in this way on the ring \(\mathbf{A}' / \mathfrak{p}'\) .

DEFINITION 3. With the notation of Definition 2, the subgroup of \(\mathcal{G}^{\mathrm{Z}}(\mathfrak{p}')\) consisting of those \(\sigma\) such that the endomorphism \(z \mapsto \sigma.z\) of \(\mathrm{A}' / \mathfrak{p}'\) is the identity is called the inertia group of \(\mathfrak{p}'\) (with respect to \(\mathcal{G}\) ) and denoted by \(\mathcal{G}^{\mathrm{T}}(\mathfrak{p}')\) (or \(\mathcal{G}^{\mathrm{T}}\) ). The ring of elements of \(\mathrm{A}'\) invariant under \(\mathcal{G}^{\mathrm{T}}(\mathfrak{p}')\) is called the inertia ring of \(\mathfrak{p}'\) (with respect to \(\mathcal{G}\) ) and is denoted by \(\mathrm{A}^{\mathrm{T}}(\mathfrak{p}')\) (or \(\mathrm{A}^{\mathrm{T}}\) ) ( \(\dagger\) ).

If A is the subring of A' consisting of the invariants of G, clearly

\[
\mathbf {A} \subset \mathbf {A} ^ {\mathbf {Z}} (\mathfrak {p} ^ {\prime}) \subset \mathbf {A} ^ {\mathbf {T}} (\mathfrak {p} ^ {\prime}) \subset \mathbf {A} ^ {\prime}.\tag{1}
\]

It follows from Definitions 2 and 3 that, for all \(\rho \in \mathcal{G}\) ,

\[
\mathcal {G} ^ {\mathrm{Z}} (\rho \cdot p ^ {\prime}) = \rho \mathcal {G} ^ {\mathrm{Z}} (p ^ {\prime}) \rho^ {- 1}, \quad \mathcal {G} ^ {\mathrm{T}} (\rho \cdot p ^ {\prime}) = \rho \mathcal {G} ^ {\mathrm{T}} (p ^ {\prime}) \rho^ {- 1}.\tag{2}
\]

If, for all \(\sigma \in \mathcal{G}^{\mathbb{Z}}(\mathfrak{p}')\) , \(\bar{\sigma}\) is the automorphism \(z \mapsto \sigma. z\) of \(\mathrm{A}' / \mathfrak{p}'\) , \(\sigma \mapsto \bar{\sigma}\) is a homomorphism (called canonical) of \(\mathcal{G}^{\mathbb{Z}}\) to the group \(\Gamma_0\) of automorphisms of

\(A^{\prime}/p^{\prime}\) leaving invariant the elements of \(A^{Z}/(p^{\prime} \cap A^{Z})\) (canonically identified with a subring of \(A^{\prime}/p^{\prime}\) ) and by definition \(\mathcal{G}^{\mathrm{T}}(p^{\prime})\) is the kernel of this canonical homomorphism; \(G^{T}\) is therefore a normal subgroup of \(G^{Z}\) . If \(k^{\prime}\) is the field of fractions of \(A^{\prime}/p^{\prime}\) , every automorphism of \(A^{\prime}/p^{\prime}\) can be extended uniquely to an automorphism of \(k^{\prime}\) , so that \(\sigma \mapsto \bar{\sigma}\) can be considered as a homomorphism from \(\mathcal{G}^{Z}(p^{\prime})\) to the group of automorphisms of \(k^{\prime}\) . Note finally that, since \(G^{T}\) is normal in \(G^{Z}\) , \(A^{T}\) is stable under \(G^{Z}\) .

LEMMA 3. Let \(\mathbf{A}'\) be a ring, \(\mathcal{G}\) a group operating on \(\mathbf{A}'\) , \(\mathbf{A}\) the ring of invariants of \(\mathcal{G}\) , \(\mathfrak{p}'\) a prime ideal of \(\mathbf{A}'\) and \(\mathbf{S}\) a multiplicative subset of \(\mathbf{A}\) not meeting \(\mathfrak{p}'\) . Then \(\mathcal{G}^{\mathrm{Z}}(\mathrm{S}^{-1}\mathfrak{p}') = \mathcal{G}^{\mathrm{Z}}(\mathfrak{p}')\) , \(\mathcal{G}^{\mathrm{T}}(\mathrm{S}^{-1}\mathfrak{p}') = \mathcal{G}^{\mathrm{T}}(\mathfrak{p}')\) and, if \(\mathcal{G}\) is locally finite, \(\mathrm{S}^{-1}\mathrm{A}^{\mathrm{Z}}(\mathfrak{p}') = \mathrm{A}^{\mathrm{Z}}(\mathrm{S}^{-1}\mathfrak{p}')\) and \(\mathrm{S}^{-1}\mathrm{A}^{\mathrm{T}}(\mathfrak{p}') = \mathrm{A}^{\mathrm{T}}(\mathrm{S}^{-1}\mathfrak{p}')\) .

As the elements of S are invariant under \(\mathcal{G}\) , clearly, if \(\sigma. p' = p'\) , also \(\sigma. (S^{-1}p') = S^{-1}p'\) . Conversely, suppose that \(\sigma \in \mathcal{G}\) is such that \(\sigma. (S^{-1}p') = S^{-1}p'\) ; then, if \(x \in p'\) , \((\sigma.x)/1 \in S^{-1}p'\) and there therefore exists \(s \in S\) such that \(s(\sigma.x) \in p'\) , whence \(\sigma.x \in p'\) since \(p'\) is prime and \(s \notin p'\) ; this proves that \(\sigma. p' \subset p'\) and it can be similarly shown that \(\sigma^{-1}. p' \subset p'\) , hence \(\sigma. p' = p'\) and \(\sigma \in \mathcal{G}^{\mathrm{Z}}(p')\) . If \(\sigma \in \mathcal{G}^{\mathrm{T}}(p')\) , then \(\sigma.x - x \in p'\) for all \(x \in A'\) , hence also, for all \(s \in S\) ,

\[
\sigma . (x / s) - (x / s) = (\sigma . x - x) / s \in \mathbf {S} ^ {- 1} \mathfrak {p} ^ {\prime}
\]

and therefore \(\sigma \in \mathcal{G}^{\mathrm{T}}(\mathrm{S}^{-1}\mathfrak{p}')\) . Conversely, suppose that \(\sigma \in \mathcal{G}^{\mathrm{T}}(\mathrm{S}^{-1}\mathfrak{p}')\) ; then, for all \(x \in \mathbf{A}'\) , \(\sigma.(x/1) - (x/1) \in \mathrm{S}^{-1}\mathfrak{p}'\) and therefore there exists \(s \in \mathbf{S}\) such that \(s(\sigma.x - x) \in \mathfrak{p}'\) , whence as above \(\sigma.x - x \in \mathfrak{p}'\) , which proves that \(\sigma \in \mathcal{G}^{\mathrm{T}}(\mathfrak{p}')\) . The last two assertions follow from § 1, no. 9, Proposition 23.

THEOREM 2. Let A' be a ring, G a finite group operating on A' and A the ring of invariants of G, so that A' is integral over A (§ 1, no. 9, Proposition 22).

\begin{enumerate}
\def\labelenumi{(\roman{enumi})}
\item
  Given two prime ideals \(\mathfrak{p}'\) , \(\mathfrak{q}'\) of \(\mathbf{A}'\) lying above the same prime ideal \(\mathfrak{p}\) of \(\mathbf{A}\) , there exists \(\sigma \in \mathcal{G}\) such that \(\mathfrak{q}' = \sigma \cdot \mathfrak{p}'\) ; in other words, \(\mathcal{G}\) operates transitively on the set of prime ideals of \(\mathbf{A}'\) lying above \(\mathfrak{p}\) .
\item
  Let \(\mathfrak{p}'\) be a prime ideal of \(\mathbf{A}', \mathfrak{p} = \mathfrak{p}' \cap \mathbf{A}\) and \(k\) (resp. \(k'\) ) the field of fractions of \(\mathbf{A}/\mathfrak{p}\) (resp. \(\mathbf{A}'/\mathfrak{p}'\) ). Then \(k'\) is a quasi-Galois extension (*) of \(k\) and the canonical homomorphism \(\sigma \mapsto \bar{\sigma}\) from \(\mathcal{G}^{\mathrm{Z}}(\mathfrak{p}')\) to the group \(\Gamma\) of \(k\) -automorphisms of \(k'\) defines, by taking the quotient, an isomorphism of \(\mathcal{G}^{\mathrm{Z}}(\mathfrak{p}') / \mathcal{G}^{\mathrm{T}}(\mathfrak{p}')\) onto \(\Gamma\) .
\setcounter{enumi}{0}
\item
  If \(x \in \mathfrak{q}'\) , then \(\prod_{\sigma \in \mathcal{G}} \sigma. x \in \mathfrak{q}' \cap A = \mathfrak{p} \subset \mathfrak{p}'\) ; then there exists \(\sigma \in \mathcal{G}\) such that \(\sigma. x \in \mathfrak{p}'\) , that is \(x \in \sigma^{-1}. \mathfrak{q}'\) . We conclude that \(\mathfrak{q}' \subset \bigcup_{\sigma \in \mathcal{G}} \sigma. \mathfrak{p}'\) and hence (since \(\mathcal{G}\) is finite and the \(\sigma. \mathfrak{p}'\) prime) there exists \(\sigma \in \mathcal{G}\) such that \(\mathfrak{q}' \subset \sigma. \mathfrak{p}'\) (Chapter II, § 1, no. 1, Proposition 2); as \(\mathfrak{q}'\) and \(\sigma. \mathfrak{p}'\) both lie above \(\mathfrak{p}, \mathfrak{q}' = \sigma. \mathfrak{p}'\) (no. 1, Corollary 1 to Proposition 1).
\end{enumerate}

\iffalse
\[
k ^ {\prime}
\]

\[
k,
\]

\[
\bar {x} \in A ^ {\prime} / p ^ {\prime}
\]

\[
k [ \mathbf {X} ]
\]

\[
\mathbf {A} ^ {\prime} / \mathfrak {p} ^ {\prime}
\]

\[
\boldsymbol {x} \in \mathbf {A} ^ {\prime}
\]

\[
\mathrm{V}, \S 6,
\]

\[
\overline {{x}};
\]
\fi

(ii) To see that \(k'\) is a quasi-Galois extension of \(k\), it suffices to prove that every element \(\bar{x}\) of \(\mathbf{A}'/\mathfrak{p}'\) is a root of a polynomial \(P\) in \(k[\mathbf{X}]\) all of whose roots are in \(\mathbf{A}'/\mathfrak{p}'\) (Algebra, Chapter V, § 6, no. 3, Corollary 3 to Proposition 9). Now, let \(x \in \mathbf{A}'\) be a representative of the class \(\bar{x}\); the polynomial

\(Q(X) = \prod_{\sigma \in \mathcal{G}} (X - \sigma.x)\) has all its coefficients in A; let \(P(X)\) be the polynomial in \((A/p)[X]\) whose coefficients are the images of those of Q under the canonical homomorphism \(\pi: A \to A/p\) . As \(\pi\) may be considered as the restriction to A of the canonical homomorphism \(\pi': A' \to A'/p'\) , it is seen that P is the product in \((A'/p')[X]\) of the linear factors \(X - \pi'(\sigma.x)\) and therefore solves the problem since \(\bar{x} = \pi'(x)\) .

Clearly, for all \(\sigma \in \mathcal{G}^{\mathbb{Z}}\) , \(\bar{\sigma}\) is a \(k\) -automorphism of \(k'\) ; it remains to verify that \(\sigma \mapsto \bar{\sigma}\) maps \(\mathcal{G}^{\mathbb{Z}}\) onto the group of all \(k\) -automorphisms of \(k'\) . Let us write \(\mathrm{S} = \mathrm{A} - \mathfrak{p}\) ; \(k\) and \(k'\) are not changed by replacing \(\mathrm{A}'\) and \(\mathfrak{p}'\) by \(\mathrm{S}^{-1}\mathrm{A}'\) and \(\mathrm{S}^{-1}\mathfrak{p}'\) respectively, by virtue of § 1, no. 9, Proposition 23 and the relation \(\mathrm{S}^{-1}\mathfrak{p}' \cap \mathrm{S}^{-1}\mathrm{A} = \mathrm{S}^{-1}(\mathrm{A} \cap \mathfrak{p}') = \mathrm{S}^{-1}\mathfrak{p}\) (Chapter II, § 2, no. 4); it follows from Lemma 3 that neither \(\mathcal{G}^{\mathbb{Z}}\) nor its operation on \(k'\) is changed; we may therefore restrict our attention to the case where \(\mathfrak{p}\) is maximal, in which case we know that so is \(\mathfrak{p}'\) (no. 1, Proposition 1) and every element of \(k'\) is therefore of the form \(\pi'(x)\) for some \(x \in \mathrm{A}'\) ; it has been seen above that such an element is a root of a polynomial in \(k[\mathrm{X}]\) of degree \(\leqslant \operatorname{Card}(\mathcal{G})\) . As every finite separable extension of \(k\) admits a primitive element (Algebra, Chapter V, § 7, no. 7, Proposition 12 and § 11, no. 4, Proposition 4), it is seen that every finite separable extension of \(k\) contained in \(k'\) is of degree \(\leqslant \operatorname{Card}(\mathcal{G})\) , whence it follows that the greatest separable extension \(k_s'\) of \(k\) contained in \(k'\) (Algebra, Chapter V, § 7, no. 6, Proposition 11) is of degree \(\leqslant \operatorname{Card}(\mathcal{G})\) (Algebra, Chapter V, § 3, no. 2, Remark 2). Let \(y \in \mathrm{A}'\) be an element such that \(\pi'(y)\) is a primitive element of \(\mathbf{k}_s'\) . The ideals \(\sigma .\mathfrak{p}'\) for \(\sigma \in \mathcal{G} - \mathcal{G}^{\mathbb{Z}}\) are maximal and distinct from \(\mathfrak{p}'\) by definition; there therefore exists \(x \in \mathrm{A}'\) such that \(x \equiv y (\text{mod. } \mathfrak{p}')\) and \(x \in \sigma^{-1}\mathfrak{p}'\) for \(\sigma \in \mathcal{G} - \mathcal{G}^{\mathbb{Z}}\) (Chapter II, § 1, no. 2, Proposition 5). This being so, let \(u\) be a \(k\) -automorphism of \(k'\) and let \(\mathrm{P}(\mathrm{X}) = \prod_{\sigma \in \mathcal{G}} (\mathrm{X} - \pi'(\sigma .x))\) ; as \(\pi'(x)\) is a root of P and \(\mathrm{P} \in k[\mathrm{X}]\) , \(u(\pi'(x))\) is also a root of P in \(k'\) and hence there exists \(\tau \in \mathcal{G}\) such that

\[
u (\pi^ {\prime} (x)) = \pi^ {\prime} (\tau . x);
\]

  but \(u(\pi'(x)) \neq 0\) and, for \(\sigma \in \mathcal{G} - \mathcal{G}^{\mathbf{Z}}\) , \(\sigma.x \in \mathfrak{p}'\) and hence \(\pi'(\sigma.x) = 0\) ; we conclude that necessarily \(\tau \in \mathcal{G}^{\mathbf{Z}}\) . But as \(u\) and \(\bar{\tau}\) have the same value for the primitive element \(\pi'(y) = \pi'(x)\) of \(k_s'\) , they coincide on \(k_s'\) and, as \(k'\) is a radicial extension of \(k_s'\) , they coincide on \(k'\) .

COROLLARY. With the hypotheses and notation of Theorem 2, let \(f_1, f_2\) be two homomorphisms of \(\mathbf{A}'\) to a field \(\mathbf{L}\) with the same restriction to \(\mathbf{A}'\) . Then there exists \(\sigma \in \mathcal{G}\) such that

\[
f _ {2} (x) = f _ {1} (\sigma . x)
\]

for all \(x \in A'\) .

Let \(p_{i}'\) be the kernel of \(f_{i}\) ( \(i = 1, 2\) ) which is a prime ideal of \(A'\) ; by hypothesis \(p_{1}' \cap A = p_{2}' \cap A\) and this intersection is a prime ideal \(p\) of \(A\) ; there therefore exists \(\tau \in \mathcal{G}\) such that \(\tau. p_{2}' = p_{1}'\) (Theorem 2 (i)); replacing \(f_{1}\) by the homomorphism \(x \mapsto f_{1}(\tau.x)\) we may then assume that \(p_{2}' = p_{1}'\) (an ideal which we shall denote by \(p'\) ). By taking the quotient we then derive from \(f_{1}\) and \(f_{2}\) two injective homomorphisms \(f_{1}', f_{2}'\) from \(A'/p'\) to \(L\) which therefore extend to two injective homomorphisms \(f_{1}''\) , \(f_{2}''\) from the field of fractions \(k'\) of \(A'/p'\) to \(L\) . As \(k'\) is a quasi-Galois extension of \(k\) , so is \(k_{1}'' = f_{1}''(k')\) and \(k_{2}'' = f_{2}''(k')\) ( \(k\) being identified with a subfield of \(L\) ) and, as there is a \(k\) -isomorphism of \(k_{1}''\) onto \(k_{2}''\) , \(k_{1}'' = k_{2}''\) (Algebra, Chapter V, § 6, Proposition 6). Thus \(f_{1}'' \circ f_{2}''\) is a \(k\) -automorphism of \(k'\) ; by Theorem 2 (ii) it is therefore of the form \(\bar{\sigma}\) , where \(\sigma \in \mathcal{G}^{Z}(p')\) . In particular, for all \(x \in A'\) the elements \(f_{2}(x)\) and \(f_{1}(\sigma.x)\) are equal.

\textbf{Remarks}\par

\begin{enumerate}
\def\labelenumi{(\arabic{enumi})}
\item
  Note that under the hypotheses of Theorem 2 \(k'\) may be infinite over \(k\) if \(k'\) is not separable over \(k\) (Exercise 9).
\item
  Clearly \(k'\) is a Galois extension of \(k\) if the field \(k\) is perfect. It is then finite over \(k\) .
\end{enumerate}

PROPOSITION 4. Let A' be a ring, G a finite group operating on A', H a subgroup of G, A and B the rings of invariants of G and H respectively and p' a prime ideal of A'; let us write p = A ∩ p' and p(B) = B ∩ p'.

\begin{enumerate}
\def\labelenumi{(\roman{enumi})}
\item
  For \(\mathcal{H}\) to be contained in the decomposition group \(\mathcal{G}^{\mathrm{Z}}(\mathfrak{p}')\) , it is necessary and sufficient that \(\mathfrak{p}'\) be the only prime ideal of A lying above \(\mathfrak{p}(\mathbf{B})\) .
\item
  If \(\mathcal{H}\) contains \(\mathcal{G}^{\mathrm{Z}}(\mathfrak{p}')\) , the following conditions are satisfied:
\end{enumerate}

\begin{enumerate}
\def\labelenumi{(\alph{enumi})}
\item
  The rings A/p and B/p(B) have the same field of fractions.
\item
  The maximal ideal of the local ring \(\mathbf{B}_{\mathfrak{p}(\mathbf{B})}\) is equal to \(\mathfrak{p}\mathbf{B}_{\mathfrak{p}(\mathbf{B})}\) .
\end{enumerate}

\begin{enumerate}
\def\labelenumi{(\roman{enumi})}
\setcounter{enumi}{2}
\item
  Suppose further that \(\mathbf{A}'\) is an integral domain and that \(\bigcap_{n\geqslant 0}\mathfrak{p}^n\mathbf{A}_{\mathfrak{p}'} = 0\) ; then conditions (a) and (b) of (ii) imply that \(\mathcal{G}^{\mathrm{Z}}(\mathfrak{p}')\) leaves invariant the elements of \(\mathbf{B}\) .

\setcounter{enumi}{0}
\item
  It follows from Theorem 2 (i) that the prime ideals of A' lying above p(B) are the ideals of the form \(\sigma. \mathfrak{p}'\) , where \(\sigma \in \mathcal{H}\) ; whence immediately (i).

\item
  We write \(S = A - p\) ; we know that the rings of invariants of \(\mathcal{G}\) and \(\mathcal{H}\) in \(S^{-1}A'\) are respectively \(S^{-1}A\) and \(S^{-1}B\) (§ 1, no. 9, Proposition 23) and \(\mathcal{G}^{\mathbb{Z}}(S^{-1}p') = \mathcal{G}^{\mathbb{Z}}(p')\) (Lemma 3); finally \(S^{-1}p(B) = S^{-1}p' \cap S^{-1}B\) (Chapter II, § 2, no. 4), the local ring of the prime ideal \(S^{-1}p(B)\) of the ring \(S^{-1}B\) is canonically isomorphic to \(B_{p(B)}\) and its residue field is isomorphic to the field of fractions of \(B/p(B)\) (Chapter II, § 2, no. 5, Proposition 11). We can therefore show (ii) restricting our attention to the case where \(p\) is maximal. To establish (a) it will be sufficient to prove that
\end{enumerate}

\[
\mathbf {B} = \mathbf {A} + \mathfrak {p} (\mathbf {B})\tag{3}
\]

for this will show that the fields A/p and B/p(B) are canonically isomorphic. By Theorem 2 there is only a finite number of prime ideals of A' lying above p and by Theorem 1 of no. 1 there is at least one prime ideal of A' lying above every prime ideal of B; this implies that there is only a finite number of prime ideals of B lying above p; let \(n_{j}\) ( \(1 \leqslant j \leqslant r\) ) denote those of these ideals which are \(\neq p(B)\) . Let x be an element of B; as the ideals \(p(B)\) and \(n_{j}\) are maximal (no. 1, Proposition 1), there exists \(y \in B\) such that \(y \equiv x \pmod{p(B)}\) and \(y \in n_{j}\) for \(1 \leqslant j \leqslant r\) (Chapter II, § 1, no. 2, Proposition 5). Let \(y_{1} = y, y_{2}, \ldots, y_{q}\) be distinct elements of the orbit of y under G; clearly

\[
z = y _ {1} + y _ {2} + \dots + y _ {q} \in \mathrm{A},
\]

and to establish (3) it will be sufficient to show that \(y_i \in \mathfrak{p}'\) for \(i \geqslant 2\) , for then we shall deduce that \(z - y \in \mathfrak{p}' \cap B = \mathfrak{p}(B)\) , whence \(x \in A + \mathfrak{p}(B)\) since \(x \equiv y \pmod{\mathfrak{p}(B)}\) . Then let \(i \geqslant 2\) and \(\sigma \in \mathcal{G}\) such that \(\sigma.y = y_i\) ; we show that \(\sigma^{-1}. \mathfrak{p}'\) does not lie above \(\mathfrak{p}(B)\) . For otherwise there would exist \(\tau \in \mathcal{H}\) such that \(\sigma^{-1}. \mathfrak{p}' = \tau. \mathfrak{p}'\) (Theorem 2 (i)), whence \((\tau^{-1}\sigma^{-1}). \mathfrak{p}' = \mathfrak{p}'\) , in other words \(\tau^{-1}\sigma^{-1} \in \mathcal{G}^{Z} \subset \mathcal{H}\) by hypothesis, whence \(\sigma \in \mathcal{H}\) ; but as \(y \in B\) and \(\sigma.y \neq y\) , this is absurd. We conclude that \(\sigma^{-1}. \mathfrak{p}'\) lies above one of the ideals \(n_j\) and, as \(y \in n_j\) by construction, certainly \(y \in \sigma^{-1}. \mathfrak{p}'\) or \(y_i = \sigma.y \in \mathfrak{p}'\) .

To prove (b) it will suffice to establish that \(\mathfrak{p}(\mathbf{B})\) is contained in the saturation \(q\) of the ideal \(\mathfrak{pB}\) with respect to \(\mathfrak{p}(\mathbf{B})\) (Chapter II, § 2, no. 4, Proposition 10); as \(\mathfrak{p}(\mathbf{B})\) is contained in none of the \(\mathfrak{n}_j\) ( \(1 \leqslant j \leqslant r\) ), it will suffice even to prove that

\[
\mathfrak {p} (\mathrm{B}) \subset \mathfrak {q} \cup \mathfrak {n} _ {1} \cup \dots \cup \mathfrak {n} _ {r}\tag{4}
\]

by Chapter II, §1, no. 1, Proposition 2. For this, we consider an element \(u \in \mathfrak{p}(\mathrm{B})\) belonging to none of the \(\mathfrak{n}_j\) ( \(1 \leqslant j \leqslant r\) ) (Chapter II, §1, no. 1, Proposition 2); let \(u_1 = u, u_2, \ldots, u_m\) be the distinct elements of the orbit of \(u\) under \(\mathcal{G}\) ; we write \(w = u_1u_2\ldots u_m, v = u_2\ldots u_m\) ; clearly \(w \in \mathrm{A}\) ; on the other hand, if \(\tau \in \mathcal{H}\) , \(\tau .u = u\) and hence necessarily \(\tau .u_{i} \neq u\) for \(i \geqslant 2\) , which shows that \(\tau .v = v\) and hence \(v \in B\) . It can be shown as in the proof of (a) that, if \(\sigma \in \mathcal{G}\) is such that \(\sigma .u = u_{i}\) where \(i \geqslant 2\) , \(\sigma^{-1}.p'\) lies above one of the \(n_{j}\) and, as \(u \notin n_{j}\) , also \(u \notin \sigma^{-1}.p'\) , in other words \(u_{i} \notin p'\) . We conclude that \(v \notin p'\) and therefore \(v \notin p(B)\) . On the other hand clearly \(w \in p' \cap A = p\) and the relation \(w = uv\) shows that \(u\) is in the saturation of \(pB\) with respect to \(p(B)\) and hence establishes (4).

\begin{enumerate}
\def\labelenumi{(\roman{enumi})}
\setcounter{enumi}{2}
\tightlist
\item
  Suppose that \(A'\) is an integral domain, that \(\bigcap_{n\geq0}p^{n}A_{p'}^{\prime}=0\) and that conditions (a) and (b) of (ii) hold. With the same notation as in (ii), clearly \(S^{-1}A'\) is an integral domain and \(S^{-1}A_{p'}^{\prime}=A_{p'}^{\prime}\) ; it is therefore possible to replace \(A'\) and \(p'\) by \(S^{-1}A'\) and \(S^{-1}p'\) , in other words, suppose also that the ideal p is maximal. The hypotheses (a) and (b) then imply that
\end{enumerate}

\[
\mathrm{B} _ {\mathfrak {p} (\mathrm{B})} = \mathrm{A} + \mathfrak {p} \mathrm{B} _ {\mathfrak {p} (\mathrm{B})}\tag{5}
\]

By induction on \(n\) we deduce that \(B_{p(B)} = A + p^n B_{p(B)}\) for all \(n > 0\) . Then let \(\sigma\) be an element of \(\mathcal{G}^{\mathbb{Z}}\) and \(x\) be an element of \(B\) . For all \(n > 0\) , there exists \(a_n \in A\) such that \(x - a_n \in p^n B_{p(B)} \subset p^n A_p'\) ; as \(\sigma .a_n = a_n\) and \(\sigma .p' = p'\) , we deduce that \(\sigma .x - x \in p^n A_p'\) . Since this relation holds for all \(n\) , we conclude from the hypothesis that \(\sigma .x = x\) .

\textbf{Remarks}\par

\begin{enumerate}
\def\labelenumi{(\arabic{enumi})}
\setcounter{enumi}{2}
\item
  If \(\mathbf{A}'\) is an integral domain and Noetherian, the condition \(\bigcap_{n\geqslant 1}\mathfrak{p}^n\mathrm{A}_{\mathfrak{p}'}' = 0\) always holds (Chapter III, §3, no. 2, Corollary to Proposition 5). It can be shown that this condition is also satisfied if \(\mathbf{A}'\) is assumed to be an integral domain and \(\mathbf{A}\) to be Noetherian.
\item
  If p is not a maximal ideal of A relation (3) does not necessarily hold under the hypotheses of (ii) and therefore A/p and B/p(B) are not necessarily isomorphic even if we take \(H = G^{Z}\) , whence \(B = A^{Z}\) (Exercise 10).
\end{enumerate}

COROLLARY 1. Under the hypotheses of Theorem 2 the rings A/p and \(\mathbf{A}^{\mathbf{Z}} / (\mathfrak{p}'\cap \mathbf{A}^{\mathbf{Z}})\) have the same field of fractions and the maximal ideal of the local ring \((\mathbf{A}^{\mathbf{Z}})_{\mathfrak{p}'\cap \mathbf{A}^{\mathbf{Z}}}\) is generated by p.

COROLLARY 2. Let A' be an integral domain, G a finite group operating on A', A the ring of invariants of G and p' a prime ideal of A'; let K, \(K^{Z}\) and K' be the fields of fractions of A, \(A^{Z}\) and A' respectively. Then K' is a Galois extension of K and the subfields L of K' containing K and such that p' is the only prime ideal of A' lying above the ideal p' ∩ L of A' ∩ L are just those which contain \(K^{Z}\) .

\(\mathcal{G}\) operates on \(\mathbf{K}'\) and \(\mathbf{K}\) is the field of invariants of \(\mathcal{G}\) in \(\mathbf{K}'\) (§ 1, no. 9, Proposition 23 applied to \(S = A - \{0\}\) ) and similarly \(\mathbf{K}^{\mathbf{Z}}\) is the field of invariants of \(\mathcal{G}^{\mathbf{Z}}\) ; by definition \(\mathbf{K}'\) is therefore a Galois extension of \(\mathbf{K}\) . If \(\mathcal{H}\) is the subgroup of \(\mathcal{G}\) consisting of those \(\sigma \in \mathcal{G}\) leaving invariant the elements of L, to say that L contains \(K^{\mathbb{Z}}\) means that \(\mathcal{H}\) is contained in \(\mathcal{G}^{\mathbb{Z}}\) (Algebra, Chapter V, § 10, no. 5, Theorem 3) and, as L is the field of invariants of \(\mathcal{H}\) in \(K'\) , \(A' \cap L\) is the ring of invariants of \(\mathcal{H}\) in \(A'\) ; the second assertion then follows from Proposition 4 (i).

DEFINITION 4. With the hypotheses and notation of Corollary 2 to Proposition 4, a prime ideal p of A is said to decompose completely in K' if the number of prime ideals of A' lying above p is equal to [K':K].

It amounts to the same to say that, for a prime ideal \(p'\) of A' lying over p, the subgroup \(\mathcal{G}^{Z}(p')\) is equal to the subgroup M leaving invariant all the elements of A', or that \(\mathrm{A}^{\mathbf{Z}}(\mathfrak{p}') = \mathrm{A}'\) , or that G/M operates faithfully on the set of prime ideals of A' lying above p.

COROLLARY 3. Let A' be an integral domain, G a finite commutative group operating on A', A the ring of invariants of G, p a prime ideal of A and K and K' the fields of fractions of A and A' respectively. Then the prime ideals of A' lying above p all have the same decomposition ring \(A^{Z}\) and the field of fractions \(K^{Z}\) of \(A^{Z}\) is the greatest intermediate field between K and K' in which p decomposes completely.

If \(\mathfrak{p}'\) is a prime ideal of A' lying above \(\mathfrak{p}\) , then \(\mathcal{G}^{\mathbf{Z}}(\sigma. \mathfrak{p}') = \mathcal{G}^{\mathbf{Z}}(\mathfrak{p}')\) since \(\mathcal{G}\) is commutative (formula (2)), hence (Theorem 2 (i)) all the prime ideals of A' lying above \(\mathfrak{p}\) have the same decomposition group \(\mathcal{G}^{\mathbf{Z}}\) and therefore the same decomposition ring \(\mathrm{A}^{\mathbf{Z}}\) ; their number is \((\mathcal{G}: \mathcal{G}^{\mathbf{Z}})\) . Let L be an intermediate field between K and K' and let \(\mathcal{H}\) be the subgroup of \(\mathcal{G}\) leaving invariant the elements of L; the decomposition group of \(\mathfrak{p}'\) with respect to \(\mathcal{H}\) is \(\mathcal{G}^{\mathbf{Z}} \cap \mathcal{H}\) ; as A' \(\cap\) L is the ring of invariants of \(\mathcal{H}\) in A', the number of prime ideals of A' lying above \(\mathfrak{p}' \cap \mathrm{L}\) is \((\mathcal{H}: (\mathcal{G}^{\mathbf{Z}} \cap \mathcal{H})) = (\mathcal{H}\mathcal{G}^{\mathbf{Z}}: \mathcal{G}^{\mathbf{Z}})\) (since \(\mathcal{G}\) is commutative). The number of prime ideals of A' \(\cap\) L lying above \(\mathfrak{p}\) is therefore \((\mathcal{G}: \mathcal{H}\mathcal{G}^{\mathbf{Z}})\) . For \(\mathfrak{p}\) to decompose completely in L, it is necessary and sufficient therefore that \((\mathcal{G}: \mathcal{H}\mathcal{G}^{\mathbf{Z}}) = [\mathrm{L}: \mathrm{K}] = (\mathcal{G}: \mathcal{H})\) and, as \(\mathcal{H} \subset \mathcal{H}\mathcal{G}^{\mathbf{Z}}\) , this is equivalent to \(\mathcal{H}\mathcal{G}^{\mathbf{Z}} = \mathcal{H}\) or also to \(\mathcal{G}^{\mathbf{Z}} \subset \mathcal{H}\) and finally to \(L \subset K^{\mathbf{Z}}\) .

PROPOSITION 5. With the hypotheses and notation of Theorem 2, the field of fractions \(k^{\mathrm{T}}\) of \(\mathrm{A}^{\mathrm{T}} / (\mathfrak{p}' \cap \mathrm{A}^{\mathrm{T}})\) is equal to the greatest separable extension \(k_s'\) of \(k\) contained in \(k'\) .

As in Proposition 4 this may be reduced to the case where p is a maximal ideal of A, which implies that \(p'\) , \(p' \cap A^{Z}\) and \(p' \cap A^{T}\) are maximal in \(A'\) , \(A^{Z}\) and \(A^{T}\) respectively (no. 1, Proposition 1).

For all \(x \in \mathbf{A}'\) , the polynomial \(\mathbf{P}(\mathbf{X}) = \prod_{\sigma \in \mathcal{G}^{\mathrm{T}}} (\mathbf{X} - \sigma.x)\) has its coefficients in the inertia ring \(\mathbf{A}^{\mathrm{T}}\) and, by definition of \(\mathcal{G}^{\mathrm{T}}\) , all its roots in \(\mathbf{A}'\) are congruent mod. \(p'\) ; the polynomial \(\pi'(\mathbf{P})\) over \(\mathbf{A}^{\mathrm{T}} / (p' \cap \mathbf{A}^{\mathrm{T}})\) whose coefficients are the canonical images of those of \(\mathbf{P}\) under the homomorphism \(\pi': \mathbf{A}' \to \mathbf{A}' / p'\) therefore has all its roots in \(\mathbf{A}' / p'\) equal to the image of \(x\) , which shows that \(k'\) is a radicial extension of \(k^{\mathrm{T}}\) ; whence \(k_{s}^{\prime} \subset k^{\mathrm{T}}\) , since every element of \(k_{s}^{\prime}\) is separable over \(k\) and a fortiori over \(k^{\mathrm{T}}\) .

We know that \(k_s'\) is a Galois extension of \(k\) (Algebra, Chapter V, § 10, no. 9, Proposition 14) and it follows from Theorem 2 that its Galois group is isomorphic to \(\mathcal{G}' = \mathcal{G}^{\mathrm{Z}} / \mathcal{G}^{\mathrm{T}}\) . As \(k^{\mathrm{T}}\) is a radicial extension of \(k_s'\) , \(k^{\mathrm{T}}\) is a quasi-Galois extension of \(k\) and the separable factor of the degree of \(k^{\mathrm{T}}\) over \(k\) is \(q = (\mathcal{G}^{\mathrm{Z}}: \mathcal{G}^{\mathrm{T}})\) . It remains to see that \(k^{\mathrm{T}}\) is a separable extension of \(k\) . We have seen above that \(\mathcal{G}'\) is identified with an automorphism group of \(A^{\mathrm{T}}\) and that \(A^{\mathrm{Z}}\) is the ring of invariants of \(\mathcal{G}'\) . If \(x \in A^{\mathrm{T}}\) , the polynomial \(Q(X) = \prod_{\sigma' \in \mathcal{G}'} (X - \sigma'(x))\) therefore has its coefficients in \(A^{\mathrm{Z}}\) ; the polynomial over \(A^{\mathrm{Z}} / (p' \cap A^{\mathrm{Z}})\) whose coefficients are the images of those of \(Q\) under \(\pi'\) is of degree \(q\) and has a root \(\pi'(x) \in A^{\mathrm{T}} / (p' \cap A^{\mathrm{T}})\) . As \(A^{\mathrm{Z}} / (p' \cap A^{\mathrm{Z}}) = k\) by Proposition 4 (ii), we see that every element of \(k^{\mathrm{T}}\) is of degree \(\leqslant q\) over \(k\) .

This being so, let \(k_{1}\) be the field of invariants of the group of \(k\) -automorphisms of the quasi-Galois extension \(k^{\mathrm{T}}\) of \(k\) ; then \([k^{\mathrm{T}}:k_{1}] = q\) (Algebra, Chapter V, § 10, no. 9, Proposition 14). Let \(u\) be a primitive element of \(k^{\mathrm{T}}\) over \(k_{1}\) ; as it is of degree \(q\) over \(k_{1}\) and of degree \(\leqslant q\) over \(k\) , it is of degree \(q\) over \(k\) and its minimal polynomial over \(k_{1}\) has coefficients in \(k\) ; this shows that \(u\) is separable over \(k\) . On the other hand, for all \(v \in k_{1}\) , there exists a power \(p^{f}\) of the characteristic exponent \(p\) such that \(v^{p^{f}} \in k\) . We conclude that \(k(u - v)\) , which contains

\[
(u - v)^{p^{f}} = u^{p^{f}} - v^{p^{f}},
\]

contains \(u^{p^{f}}\) and consequently \(k(u^{p^{f}})\) . But as \(u\) is separable over \(k\) , \(k(u) = k(u^{p^{f}})\) (Algebra, Chapter V, § 3, no. 3, Proposition 4), whence \(k(u) \subset k(u - v)\) . As \(u\) is of degree \(q\) over \(k\) and \(u - v\) of degree \(\leqslant q\) , it follows that \(k(u) = k(u - v)\) , whence \(v \in k(u)\) . This shows that \(v\) is separable over \(k\) , hence \(k_1 = k\) and \(k^{\mathrm{T}}\) is separable over \(k\) .

COROLLARY. If the order of the inertia group \(\mathcal{G}^{\mathrm{T}}(\mathfrak{p}')\) is relatively prime to the characteristic exponent \(p\) of \(k\) , the field \(k'\) is a Galois extension of \(k\) .

With the notation of the proof of Proposition 5, the polynomial \(\pi'(\mathbf{P})\) has coefficients in \(k^{\mathrm{T}} = k_{s}'\) and all its roots equal to \(\pi'(x)\) ; we immediately deduce that \(\pi'(\mathbf{P})\) is a power of a minimal polynomial of \(\pi'(x)\) over \(k_{s}'\) ; but the latter has degree equal to a power of \(p\) and hence, as the degree of \(\pi'(\mathbf{P})\) is equal to the order of \(\mathcal{G}^{\mathrm{T}}\) , the hypothesis implies that \(\pi'(x) \in k_{s}'\) , in other words \(k_{s}' = k'\) .
% semantic-fragments: legacy-input-seam
\relax{}
